%% Paper skeleton for the JXUST modeling class (MIT).
%%
%% This is the ModelForge adapter for `JXUSTmodeling.cls` (vendored from
%% sikouhjw/JXUSTmodeling, MIT). The upstream class defines its own front matter
%% commands — `\biaoti` for the title and `\keyword` for the keywords — and an
%% `abstract` environment plus an `appendixx` environment, so the body is written
%% against those rather than against ctexart.
%%
%% Compile:  latexmk -xelatex -interaction=nonstopmode paper.tex
%% Needs:    JXUSTmodeling.cls in the same directory (it is).
%%
%% See `example.tex` in this directory for the upstream's full feature tour
%% (tables, TikZ figures, flow charts).
%%
\documentclass{JXUSTmodeling}

\biaoti{论文标题}
\keyword{关键词一；关键词二；关键词三}

\begin{document}

\begin{abstract}
  摘要用 300--500 字概括：针对什么问题、建立了什么模型、用什么方法求解、
  得到了哪些主要定量结论、模型有何特点。先写完全文结果，再回来完成摘要。

  \textbf{针对问题一，} 建立了\ldots 模型，采用\ldots 求解，得到\ldots。

  \textbf{针对问题二，} 应用了\ldots 方法，发现\ldots。

  \textbf{针对问题三，} 提出了\ldots，其优势在于\ldots。
\end{abstract}

\section{问题的提出}\label{sec:1}
\subsection{问题的背景}\label{sub:1.1}
\subsection{已知的条件}\label{sub:1.2}
\subsection{问题的提出}\label{sub:1.3}

\section{问题的分析}\label{sec:2}
\subsection{问题的整体分析}\label{sub:2.1}
\subsection{问题一的分析}\label{sub:2.2}
\subsection{问题二的分析}\label{sub:2.3}

\section{模型的假设}\label{sec:3}
为简化问题，本文做出以下假设：
\begin{enumerate}
  \item 假设一\ldots
  \item 假设二\ldots
\end{enumerate}

\section{符号说明}\label{sec:4}
\begin{table}[htbp]
  \centering
  \caption{符号说明表}
  \begin{tabular}{cll}
    \toprule
    符号 & 说明 & 单位 \\
    \midrule
    $m$ & 质量 & kg \\
    $t$ & 时间 & s \\
    \bottomrule
  \end{tabular}
\end{table}

\section{模型的建立与求解}\label{sec:5}
\subsection{问题一的模型建立}
给出目标函数、决策变量与约束条件，并说明模型选择的理由与局限。
\begin{equation}
  \label{eq:objective}
  \min_{x \in X} f(x) \quad \text{s.t.} \quad g_i(x) \le 0 .
\end{equation}

\subsection{问题一的模型求解}
说明算法步骤、参数取值与实现方式，固定随机种子以便复现。

\subsection{求解结果与分析}
给出求解结果；图表紧随首次引用处。

\section{模型的检验}\label{sec:6}
给出灵敏度分析、误差指标与与基线方法的对比。

\section{模型的评价及推广}\label{sec:7}
列出模型的优点、局限与可推广的方向。

\begin{thebibliography}{99}
  \bibitem{ref1} 作者. 论文名. 期刊, 年份.
\end{thebibliography}

\begin{appendixx}
  \section{问题一的求解代码}
\end{appendixx}

\end{document}
